% Fragmento para inclusão no artigo/tese.
% Requer no preâmbulo: graphicx, booktabs e url (ou hyperref), além do comando \source.

\subsection{Experimentos com a abordagem \textit{Ontology-Based}}

Os experimentos com a abordagem \textit{ontology-based} foram realizados por
meio de uma API desenvolvida com o \textit{framework} Flask. O serviço recebe um
texto, combina-o com um \textit{prompt few-shot} e permite que o modelo de
linguagem consulte informações ontológicas da Wikidata antes de produzir o
grafo de conhecimento em RDF/Turtle. O acesso à Wikidata é realizado por meio
de um servidor compatível com o \textit{Model Context Protocol} (MCP)%
\footnote{Endpoint utilizado: \url{https://wd-mcp.wmcloud.org/mcp/}.}.

Essa configuração corresponde à ablação intermediária entre os experimentos
baseados exclusivamente em \textit{prompting} e a abordagem de
\textit{hybrid pipelines}. No primeiro caso, o modelo recebe apenas as
instruções e os exemplos contidos nos \textit{prompts}. Na abordagem
\textit{ontology-based}, o modelo também pode buscar identificadores e
hierarquias de classificação na Wikidata. No fluxo híbrido, por sua vez, a
aplicação executa explicitamente etapas de extração de menções, resolução de
entidades, recuperação de declarações e identificação de relações diretas. A
Tabela~\ref{tab:comparacao-ablacoes-kg} resume essa delimitação.

\begin{table}[htbp]
\centering
\caption{Delimitação das abordagens avaliadas na construção dos grafos.}
\label{tab:comparacao-ablacoes-kg}
\small
\renewcommand{\arraystretch}{1.25}
\begin{tabular}{p{3.4cm}p{3.2cm}p{3.4cm}p{4.1cm}}
\toprule
\textbf{Abordagem} &
\textbf{Uso de \textit{prompt}} &
\textbf{Acesso à Wikidata} &
\textbf{Orquestração local} \\
\midrule
\textit{Prompt-based} &
Sistema e exemplos \textit{few-shot} &
Não utiliza recuperação externa &
Composição do \textit{prompt}, reparo e validação do RDF \\
\textit{Ontology-based} &
Sistema e exemplos \textit{few-shot} &
Busca de itens e hierarquia de instância/subclasse por MCP &
Execução das chamadas de ferramenta, reparo e validação do RDF \\
\textit{Hybrid pipelines} &
Sistema, extração e construção do RDF &
Entidades, declarações e relações estruturadas &
Extração, resolução, enriquecimento, construção e validação \\
\bottomrule
\end{tabular}
\end{table}

\begin{figure}[htbp]
    \centering
    \caption{Fluxo de criação de grafos de conhecimento com a abordagem
    \textit{ontology-based}.}
    \includegraphics[width=0.52\linewidth]{metodo/figuras/experimento-b.jpg}
    \label{fig:experimento-b}
    \source{da autora.}
\end{figure}

Conforme apresentado na Figura~\ref{fig:experimento-b}, o fluxo inicia-se com
uma requisição HTTP do tipo \texttt{POST} enviada à rota
\texttt{/analyze}. O corpo deve conter o campo textual \texttt{text} e pode
informar os arquivos de \textit{prompt} e o número máximo de tentativas de
geração do RDF. O controlador verifica se o corpo é um objeto JSON, se o texto
é uma cadeia não vazia e se os parâmetros opcionais possuem os tipos esperados.
O texto é normalizado pela remoção de espaços em suas extremidades.

Em seguida, o serviço carrega o \textit{prompt} de sistema e o
\textit{prompt few-shot}. O texto é inserido no marcador
\texttt{\$\{USER\_TEXT\}}; também é aceito o marcador legado
\texttt{\$\{Text\_TEXT\}}. Quando nenhum dos dois está presente, a aplicação
acrescenta o conteúdo no formato de uma interação entre usuário e assistente.

Antes da geração, o cliente MCP inicializa uma sessão por HTTP, envia a
notificação de inicialização e consulta o método \texttt{tools/list}. As
ferramentas descobertas são filtradas por uma lista de permissão configurável.
No perfil experimental padrão, apenas duas operações são expostas ao modelo:
\texttt{search\_items}, destinada à busca de itens e identificadores, e
\texttt{get\_instance\_and\_subclass\_hierarchy}, destinada à recuperação da
hierarquia formada pelas propriedades \textit{instance of} e
\textit{subclass of}. Ferramentas capazes de recuperar declarações arbitrárias
ou executar consultas SPARQL não são disponibilizadas nessa ablação.

Os esquemas JSON das ferramentas permitidas são convertidos para o formato de
funções aceito pela rota \texttt{/api/chat} do Ollama. A requisição contém o
modelo configurado, o histórico de mensagens, as ferramentas e as opções de
geração, com \texttt{stream} definido como \texttt{false}. O modelo utilizado
deve oferecer suporte a chamadas estruturadas de ferramenta; por esse motivo,
o perfil padrão emprega \texttt{llama3.1:8b}.

Quando o modelo solicita uma ferramenta, a aplicação interpreta o nome e os
argumentos da função, verifica se a operação pertence à lista de permissão e
executa uma chamada MCP do tipo \texttt{tools/call}. O idioma configurado é
incluído automaticamente quando o esquema da ferramenta oferece o argumento
\texttt{lang}. O conteúdo retornado pela Wikidata é acrescentado ao histórico
como uma mensagem com papel \texttt{tool}, após a qual o histórico completo é
novamente enviado ao modelo. Várias chamadas devolvidas na mesma resposta são
executadas sequencialmente.

O ciclo entre o modelo e as ferramentas continua até que seja produzida uma
resposta textual sem novas chamadas ou até que seja atingido o limite definido
por \texttt{MAX\_TOOL\_ROUNDS}. Com
\texttt{REQUIRE\_WIKIDATA\_MCP=true}, uma resposta final produzida antes de
qualquer consulta é rejeitada provisoriamente e o modelo recebe uma instrução
adicional para utilizar a Wikidata. Caso uma segunda resposta seja produzida
sem chamadas de ferramenta, o processamento é encerrado com erro. Essa regra
evita que uma execução configurada como \textit{ontology-based} seja
silenciosamente reduzida ao experimento baseado somente em \textit{prompt}.

O requisito implementado verifica a ocorrência de ao menos uma chamada MCP. A
seleção específica entre busca e consulta da hierarquia permanece orientada
pelos \textit{prompts}. Portanto, o modelo é instruído a pesquisar as entidades
reconhecíveis e consultar a hierarquia dos identificadores selecionados quando
a classificação for relevante, mas o serviço não impõe uma sequência fixa de
ferramentas para todos os textos.

\subsubsection{Definição dos \textit{prompts}}

Foram utilizados dois arquivos com funções complementares: o \textit{prompt}
de sistema \texttt{prompt/system/knowledge\_graph.txt} e o \textit{prompt}
\textit{few-shot}
\texttt{prompt/prompts/ontology-few-shot.txt}. Essa organização preserva a
estrutura do experimento \textit{prompt-based}, permitindo que o efeito da
ancoragem ontológica seja analisado sem introduzir um \textit{prompt} separado
para extração de menções.

O \textit{prompt} de sistema define o modelo como um assistente de dados
estruturados fundamentado em ontologias. Ele determina que entidades do mundo
real sejam pesquisadas com \texttt{search\_items} e que identificadores
selecionados tenham sua hierarquia inspecionada quando necessário. Os
identificadores retornados pelo MCP durante a requisição corrente constituem a
única fonte autorizada de \textit{Q-IDs}; identificadores recordados pelo modelo
ou apenas exibidos nos exemplos não devem ser reutilizados sem confirmação.

O \textit{prompt few-shot} mantém as instruções de serialização, os prefixos e
os três exemplos de transformação utilizados no experimento baseado em
\textit{prompting}. Entretanto, explicita que os \textit{Q-IDs} presentes nos
exemplos demonstram somente o formato esperado. Também separa dois tipos de
evidência: a hierarquia obtida na Wikidata pode fundamentar relações de
classificação, enquanto as demais relações representadas no grafo devem ser
apoiadas pelo texto de entrada. Assim, uma declaração disponível na Wikidata,
mas ausente do texto, não deve ser adicionada como fato narrado pelo documento.

Os \textit{prompts} exigem que a resposta final contenha exclusivamente
RDF/Turtle. Os recursos utilizados como sujeito ou objeto devem possuir
\texttt{rdfs:label} com marcação de idioma. Recursos confirmados pelo MCP são
representados com o prefixo \texttt{wd:}; conceitos não resolvidos recebem
identificadores locais no \textit{namespace} \texttt{kg:}. Para classificações,
o predicado preferencial é \texttt{kg:is}. Também são priorizadas relações
entre recursos rotulados, pois elas permitem que consultas SPARQL percorram o
grafo até as entidades correspondentes às respostas esperadas.

\subsubsection{Validação e correção do RDF/Turtle}

Quando o modelo retorna uma resposta textual, o serviço identifica o trecho
correspondente ao RDF e remove delimitadores de blocos de código ou conteúdo
anterior à primeira marca reconhecível de Turtle. Em seguida, são gerados
candidatos de reparo para erros recorrentes, como aspas tipográficas, aspas
duplicadas, ausência de ponto final, uso do prefixo \texttt{xsd:} sem declaração
e identificadores da Wikidata incompletos. Blocos finais interrompidos também
podem ser removidos quando já existe uma declaração completa.

Cada candidato é analisado por meio de
\texttt{rdflib.Graph.parse(format="turtle")}. Para ser aceito, o resultado deve
ser sintaticamente válido, conter pelo menos uma tripla e possuir ao menos um
predicado diferente de \texttt{rdfs:label}. Na última tentativa, a aplicação
também pode preservar somente as declarações que forem analisáveis de forma
independente.

Se nenhum candidato for válido e ainda houver tentativas disponíveis, o erro
do analisador e a saída anterior são acrescentados ao mesmo histórico de
conversa. O modelo recebe então uma solicitação de correção. As evidências MCP
já presentes no histórico são reutilizadas, de modo que uma consulta bem-sucedida
à Wikidata não precisa ser repetida apenas porque o RDF inicial estava
malformado. O parâmetro \texttt{max\_rdf\_attempts} é limitado ao intervalo de
uma a três tentativas. O esgotamento das tentativas resulta em HTTP
\texttt{422}.

Além do texto e do RDF validado, a resposta bem-sucedida contém os nomes dos
\textit{prompts}, a atribuição \texttt{Source: Wikidata MCP} e a lista
\texttt{mcp\_calls}. Essa lista registra, em ordem, o nome, os argumentos e uma
versão limitada do resultado de cada ferramenta. O resultado conservado na
resposta possui no máximo 6.000 caracteres, embora o conteúdo completo seja
fornecido ao modelo durante a execução. As chamadas ao Ollama também são
registradas em CSV com o histórico de mensagens, a resposta textual e as
chamadas estruturadas de ferramenta.

\subsubsection{Configuração do experimento}

A Tabela~\ref{tab:config-ontology-based} apresenta os principais parâmetros do
perfil padrão da abordagem \textit{ontology-based}.

\begin{table}[htbp]
\centering
\caption{Configurações do experimento com a abordagem \textit{ontology-based}.}
\label{tab:config-ontology-based}
\small
\renewcommand{\arraystretch}{1.25}
\begin{tabular}{p{4.2cm}p{2.2cm}p{6.2cm}p{2.4cm}}
\toprule
\textbf{Parâmetro} & \textbf{Tipo} & \textbf{Descrição} & \textbf{Valor padrão} \\
\midrule
\texttt{OLLAMA\_MODEL} & texto & Modelo com suporte a ferramentas & \texttt{llama3.1:8b} \\
\texttt{stream} & booleano & Desabilita resposta incremental & \texttt{false} \\
\texttt{OLLAMA\_TEMPERATURE} & real & Controla a aleatoriedade & \texttt{0} no Compose \\
\texttt{OLLAMA\_NUM\_PREDICT} & inteiro & Limite de \textit{tokens} gerados & \texttt{1536} no Compose \\
\texttt{WIKIDATA\_MCP\_TOOLS} & lista & Ferramentas expostas ao modelo & busca e hierarquia \\
\texttt{WIKIDATA\_TIMEOUT\_SECONDS} & real & Tempo máximo por requisição MCP & \texttt{60} \\
\texttt{REQUIRE\_WIKIDATA\_MCP} & booleano & Exige ao menos uma chamada MCP & \texttt{true} \\
\texttt{MAX\_TOOL\_ROUNDS} & inteiro & Limite de ciclos modelo--ferramenta & \texttt{8} \\
\texttt{max\_rdf\_attempts} & inteiro & Tentativas de geração/correção do RDF & \texttt{3} \\
\bottomrule
\end{tabular}
\end{table}

Os parâmetros de amostragem do Ollama são opcionais e somente são incluídos na
requisição quando possuem valor. A configuração do Docker Compose utiliza
temperatura zero e limite de 1.536 \textit{tokens}, enquanto uma execução local
sem essas variáveis utiliza os valores internos do modelo para os parâmetros
não especificados.

\subsubsection{Considerações sobre a reprodutibilidade}

A abordagem depende de três fontes de variabilidade. A primeira corresponde à
geração do modelo e às decisões sobre quais ferramentas chamar, quais
candidatos selecionar e quando encerrar o ciclo. A segunda decorre da Wikidata,
cujo conteúdo, ordenação de resultados e hierarquias podem mudar após a coleta.
A terceira está relacionada aos reparos locais, pois saídas inicialmente
distintas podem ser transformadas em diferentes candidatos válidos.

Para favorecer a rastreabilidade, recomenda-se registrar a versão do modelo,
as opções de geração, a data da execução, o endereço do servidor MCP e a
resposta completa do experimento. O campo \texttt{mcp\_calls} permite verificar
quais ferramentas foram executadas, mas seus resultados são truncados na
resposta pública; quando a reprodução integral do conteúdo externo for
necessária, deve-se conservar também o tráfego ou uma cópia completa das
evidências utilizadas.

Em síntese, o experimento \textit{ontology-based} mantém a geração orientada por
\textit{few-shot prompting}, mas restringe o uso de identificadores e classes a
evidências obtidas por meio do MCP da Wikidata. A aplicação coordena as chamadas
de ferramenta e valida o RDF, sem executar o processo explícito de extração e
enriquecimento relacional característico do \textit{hybrid pipeline}. Dessa
forma, a abordagem permite avaliar o efeito da ancoragem ontológica mantendo
uma separação operacional clara entre os três níveis de ablação.
